\uuid{nEkp}
\titre{Nombre de paramètres d’un réseau de neurones}
\niveau{L3}
\module{Analyse de données}
\chapitre{Réseaux de neurones}
\sousChapitre{Architecture, régression multiple}
\theme{Nombre de paramètres, régression multiple, couche cachée}
\auteur{Maxime Nguyen}
\datecreate{2026-03-19}
\organisation{AMSCC}
\difficulte{1}

\contenu{
	\texte{
		On considère deux réseaux de neurones ayant une seule couche cachée contenant $h$ neurones et une unique sortie.
		
		Le premier réseau reçoit une seule variable d'entrée : son architecture est $1 \to h \to 1$.
		
		Le second réseau reçoit trois variables d'entrée : son architecture est $3 \to h \to 1$.
		
		L'objectif est de comprendre comment le nombre total de paramètres dépend du nombre de variables d'entrée, lorsque le nombre de neurones cachés reste fixé.
	}
	\begin{enumerate}
		
		\item
		\question{Exprimer en fonction de $h$ le nombre de paramètres du réseau $3 \to h \to 1$. Détailler par couche.}
		\reponse{
			\begin{itemize}
				\item Couche d'entrée vers couche cachée : $3h$ poids et $h$ biais, soit $4h$ paramètres.
				\item Couche cachée vers sortie : $h$ poids et $1$ biais, soit $h+1$ paramètres.
			\end{itemize}
			Le nombre total de paramètres est donc
			$$
			4h + (h+1) = \mathbf{5h+1}.
			$$
		}
		
		\item
		\question{Exprimer en fonction de $h$ le nombre de paramètres du réseau $1 \to h \to 1$.}
		\reponse{
			\begin{itemize}
				\item Couche d'entrée vers couche cachée : $h$ poids et $h$ biais, soit $2h$ paramètres.
				\item Couche cachée vers sortie : $h$ poids et $1$ biais, soit $h+1$ paramètres.
			\end{itemize}
			Le nombre total de paramètres est donc
			$$
			2h + (h+1) = \mathbf{3h+1}.
			$$
		}
		
		\item
		\question{Combien de paramètres supplémentaires le réseau $3 \to h \to 1$ possède-t-il par rapport au réseau $1 \to h \to 1$ ? Interpréter ce résultat.}
		\reponse{
			La différence vaut
			$$
			(5h+1) - (3h+1) = \mathbf{2h}.
			$$
			
			Le réseau à trois entrées possède donc $\mathbf{2h}$ paramètres supplémentaires.
			
			Cela s'explique par le fait que, par rapport au réseau à une seule entrée, on ajoute $2$ variables d'entrée supplémentaires, chacune étant connectée aux $h$ neurones de la couche cachée. Cela ajoute donc $2h$ poids.
		}
		
		\item
		\question{Application numérique : retrouver le résultat lorsque $h=8$.}
		\reponse{
			Si $h=8$, alors :
			$$
			5h+1 = 5 \times 8 + 1 = 41
			$$
			pour le réseau $3 \to 8 \to 1$, et
			$$
			3h+1 = 3 \times 8 + 1 = 25
			$$
			pour le réseau $1 \to 8 \to 1$.
			
			La différence est donc
			$$
			41 - 25 = 16.
			$$
		}
		
	\end{enumerate}
}